\section{Related Work}

\subsection{Knowledge Graph Construction in Cybersecurity}

Several efforts have focused on constructing Cybersecurity Knowledge Graphs (CSKGs) to integrate and structure diverse threat intelligence sources. One of the most established initiatives is the SEPSES CSKG, which unifies multiple structured cybersecurity datasets—including CVE, CAPEC, CWE, ATT\&CK, and NVD—into a consolidated semantic graph \cite{Kiesling2019SEPSES}. The SEPSES framework demonstrates the value of applying ETL workflows and semantic integration for supporting tasks such as vulnerability assessment, attack pattern exploration, and cross-dataset reasoning. However, its reliance on structured and semi-structured data limits its ability to capture the rich and dynamic information contained within unstructured Cyber Threat Intelligence (CTI) reports.

Beyond SEPSES, several ontologies contribute foundational modeling constructs for cybersecurity. The Unified Cybersecurity Ontology (UCO) provides a harmonized and interoperable vocabulary that bridges heterogeneous cybersecurity schemas and supports the representation of events, indicators of compromise (IoCs), threat actors, and system behaviors \cite{UCOCommunity}. Meanwhile, the MALOnt ontology focuses specifically on malware behavior and system interactions, offering a structured framework for modeling how malicious software propagates and affects digital infrastructures \cite{Rastogi2020MALOnt}. These ontologies, while comprehensive, assume standardized and well-formatted inputs, and therefore do not directly address the extraction or integration of knowledge from unstructured text.

More recent work extends CSKG construction into specialized domains. The ICS-SEC Knowledge Graph applies semantic integration techniques to the industrial control systems (ICS) security domain, combining ICS-specific ontologies, structured IoC feeds, and mappings to standards such as STIX \cite{Kurniawan2024ICSSEC}. Although ICS-SEC expands the applicability of CSKGs to operational technology (OT) environments, it—like SEPSES—remains reliant on structured intelligence sources and does not include automated extraction from narrative CTI artifacts such as analyst blogposts or incident reports.

Collectively, these initiatives highlight significant progress in structured-data CSKG development while revealing a persistent gap: the lack of mechanisms for incorporating knowledge from unstructured CTI sources. This limitation motivates the need for automated extraction and ontology-driven integration techniques to enrich CSKGs with information derived from natural-language threat intelligence.

\subsection{Unstructured Threat Intelligence Extraction}

A substantial portion of actionable cyber threat intelligence is expressed in unstructured textual formats, including incident reports, security advisories, malware analyses, and technical blogposts. Extracting structured knowledge from these free-text sources has therefore become an important research direction within cybersecurity natural language processing (NLP) and information extraction (IE).

Early approaches primarily relied on rule-based extraction, using handcrafted patterns or regular expressions to identify basic indicators of compromise (IoCs) such as IP addresses, URLs, file hashes, and registry keys. While effective for simple artifacts, such methods lack generalizability and are unable to capture the deeper behavioral and relational information present in narrative CTI.

Later studies introduced traditional machine learning models, particularly named entity recognition (NER) models trained on limited cybersecurity corpora, to extract domain-specific entities such as malware names, CVE identifiers, threat actors, and attack techniques. Although these approaches improved robustness, they remained constrained by the scarcity of annotated CTI datasets and the high linguistic variability of cybersecurity text.

Recent advancements leverage deep neural networks and Large Language Models (LLMs) to extract both entities and relations from unstructured CTI. These approaches enable identification of richer semantic concepts such as attacker behavior, exploitation chains, and tactics, techniques, and procedures (TTPs). For example, ontology-guided extraction frameworks such as OntoLogX employ LLMs and semantic alignment mechanisms to transform cybersecurity log data and free-text narratives into structured knowledge graph triples \cite{Cotti2025OntoLogX}. Some pipelines further incorporate mapping to standardized cybersecurity vocabularies such as STIX to normalize extracted intelligence.

Despite these developments, end-to-end pipelines that incorporate unstructured CTI extraction into full CSKG construction remain rare. Challenges persist due to inconsistent terminology, implicit relationships, narrative reporting structures, and limited labeled data. As a result, most existing CSKGs remain grounded in structured datasets. Addressing these limitations, this project aims to develop a pipeline capable of extracting, normalizing, and semantically integrating unstructured CTI into a CSKG that complements existing structured resources such as SEPSES and ICS-SEC.
